% Generated draft: requires mathematical and source review.
\documentclass[11pt]{article}
\usepackage[margin=25mm]{geometry}
\usepackage{fontspec}
\setmainfont{Latin Modern Roman}
\usepackage{amsmath,amssymb,mathtools,unicode-math}
\setmathfont{Latin Modern Math}
\usepackage{enumitem,needspace,xurl,hyperref}
\setlength{\parindent}{0pt}
\setlength{\parskip}{5pt}
\setlength{\emergencystretch}{3em}
\newcommand{\N}{\mathbb{N}}
\newcommand{\Z}{\mathbb{Z}}
\newcommand{\Q}{\mathbb{Q}}
\newcommand{\R}{\mathbb{R}}
\newcommand{\C}{\mathbb{C}}
\newcommand{\F}{\mathbb{F}}
\newcommand{\A}{\mathbb{A}}
\newcommand{\PP}{\mathbb{P}}
\newcommand{\E}{\mathbb{E}}
\newcommand{\eps}{\varepsilon}
\newcommand{\dd}{\,\mathrm{d}}
\begin{document}
\section*{UnsolvedMath ID: 10600038}
\textbf{Source problem code:} catalogue entry
\subsection{Definitions and mathematical statement}
A virtual link diagram has classical crossings, carrying over/under data, and virtual crossings. Its virtual isotopy class is generated by the generalized Reidemeister moves, including the virtual detour moves. The source asks for a move description within a restricted class of such diagrams; it does not specify whether only one-component diagrams are intended.

The phrase ``leading term equal to zero'' refers to cancellation at an a priori highest possible degree, rather than to the actual highest nonzero coefficient of a polynomial. For example, in the usual normalization for a nonempty diagram,
\[
 \langle D\rangle=\sum_s A^{a(s)-b(s)}(-A^2-A^{-2})^{|s|-1},
\]
where the states smooth the classical crossings, $a(s)$ and $b(s)$ count the two smoothing types, and $|s|$ counts the resulting state circles. If $c(D)$ is the number of classical crossings and $s_A$ is the all-$A$ state, the customary candidate upper degree is
\[
 M(D)=c(D)+2(|s_A|-1).
\]
Under this convention, write
\[
 \mathcal C=\{D:[A^{M(D)}]\langle D\rangle\ne0\}.
\]
The source does not explicitly give this normalization or degree formula; the intended highest-possible-degree convention must be checked against the survey before treating this displayed definition as definitive. In particular, the condition is stronger than merely requiring a nonzero bracket polynomial.

\textbf{Problem.} Classify moves, each a composition of generalized Reidemeister moves, that establish virtual isotopy equivalence within the class of diagrams having nonzero highest-possible-degree term. A complete move description should connect any two diagrams $D,D'\in\mathcal C$ representing the same virtual isotopy class by a sequence of the specified moves with diagrams in $\mathcal C$ between the specified moves. The source does not state whether every elementary intermediate diagram inside a composite move must also belong to $\mathcal C$, nor whether a finite list of composite move types is required. These distinctions are unresolved parts of the formulation, rather than additional assumptions.
\subsection{Short English statement}
Describe the combinations of virtual Reidemeister moves needed to relate equivalent diagrams whose Kauffman brackets retain their expected top-degree term. The expected degree convention and the admissibility of intermediate diagrams need clarification.
\subsection{Sources}
\begin{itemize}
\item Local UnsolvedMath registry, permanent ID 10600038, source code catalogue entry, including its complete statement, provenance, and research summary.
\item \emph{Unsolved Problems in Virtual Knot Theory and Combinatorial Knot Theory}, revised survey, problem-list item 38; \url{https://arxiv.org/abs/1409.2823}. The registry records source-TeX extraction and access on 2026-07-29. No individual proposer or proposal year is supplied.
\item The registry summary mentions ``Local transformations and functorial maps'' (2023) with the literal identifier arXiv:2301.xxxx, and unspecified parity-based bracket studies. The identifier is a placeholder, not a usable citation; the studies are not individually identified.
\end{itemize}
\end{document}
